\cotefile{79}
\scanpage{220}
qui \unsure{n'est autre|n'est-autre} que l'iso. \inserted{du type} (76), pour le système
dual du système envisagé \ldots

\deleted{\illegible{1} Il \unsure{reste donc qu'ainsi|reste donc que si l'on}}

\uline{Transport de structure}. Soit $g$ un autom. de
$(S,S^\vee,R)$, et soit $\varphi:S\longrightarrow S^\vee$ une jauge (il n'est
pas question d'un $Z$, pour l'instant). Alors \deleted{\illegible{1}}
$\varphi$ définit une structure de graphe $A_\varphi$, et on
a par transport de structure
\[
\left\{\begin{aligned}
g(\varphi)&\overset{\mathrm{déf}}=g_{S^\vee}\circ\varphi\circ g_S^{-1} &&\text{d'où :}\\
g_S(A_\varphi)&=A_{g(\varphi)}
\end{aligned}\right. \tag*{(87)}
\]

Quand on dispose d'un $Z$ ss-groupe comm. de
$\mathrm{Aut}(\Gamma=(S,S^\vee,\mathcal R))$, \inserted{$\varphi$ est une $g$-jauge, par ex. quand}
\inserted{et $g\in Z$, $\varphi_0\in\Pi_Z$, alors}
\deleted{du centralisateur $G$ dans $\mathrm{Aut}(\Gamma)$,}
\deleted{alors si $g\in G$ est une}
\deleted{alors l'ens. de $Z$-jauges $\varphi$ -- i.e. satisfaisant}
\deleted{$\lambda_{S^\vee}\varphi=\varphi\lambda_S^{-1}\quad\forall\lambda\in Z$, \unsure{déduction} \unsure{vide}, par une $G$-torsion via}
\[
 T_g\varphi_0=g_{S^\vee}\varphi_0=\varphi_0g_S^{-1} \tag*{(88)}
\]

\deleted{(NB que quelque soit $g$, si $Z$-jauges fixée, et $g\in\mathrm{Aut}(\Gamma)$,)}
\[
 \deleted{g_{S^\vee}\varphi=\varphi g_S^{-1}\quad\Longleftrightarrow\quad
 g\in\mathrm{Cent}_{\mathrm{Aut}(\Gamma)}(Z)\;(=G)} \tag*{\deleted{(89)}}
\]

Donc, on a alors, \deleted{pour $g\in G$}
\[
 \forall g\in\deleted{\unsure{G}}Z,\quad
 \varphi_0\in\underbrace{\Pi_Z}_{\substack{\text{ens. des}\\ Z\text{-jauges}}}
 \qquad \underbrace{g(\varphi)}_{\text{cf (87)}}=\deleted{\unsure{g}}T_{g^2}\varphi
 \tag*{(90)}
\]

Plus généralement (pour mémoire) si $\pi$ est un sous-$Z$-tors. de $\Pi_Z$,
\[
\begin{aligned}
&\text{Si }g\in G=\mathrm{Cent}_{\mathrm{Aut}(\Gamma)}(Z),\text{ alors conditions équivalentes}\\
&\left\{\begin{aligned}
\text{a)}\quad g_{S^\vee}\varphi&=\varphi g_S^{-1} &&\text{pour \uline{un} }\varphi\in\pi\quad(\text{disons, }\varphi_0)\\
\text{b)}\quad g_S\varphi&=\varphi g_{S^\vee}^{-1} &&\text{pour tt }\varphi\in\pi
\end{aligned}\right.\\
&\hspace{3em}(\text{NB relation qui s'écrit }\varphi(g_S)=(g_{S^\vee})^{-1}
\end{aligned}\tag*{(91)}
\]
% slanted
\margin{\uline{NB} Correspondent aux homothéties pour les structures $A_\varphi$ ($\varphi\in\pi$).
Les $g$ forment \deleted{sont \unsure{spéciales}} \unsure{un sous|un ensemble} \deleted{\unsure{stratification|stabilisateur}} (\unsure{i.e.|éqv.} sous \deleted{\unsure{Z}} s.g.) c'est pareil) sur $Z$}
et alors on \deleted{\unsure{a}} a
\[
 g(\varphi)=T_{g^2}\varphi\qquad\text{pour toute }\varphi\in\pi.
 \tag*{(92)}
\]
% archive sheet number: 219

\scanpage{221}
\mspage{113}
\deleted{4} Retour sur la condition d'isomorphie.

% slanted
\margin{(on suppose $R$ séparant en $S^\vee$)}
Soit
\[
 \deleted{\unsure{G}}\quad A=\mathrm{Aut}(\Gamma),\qquad
 \Gamma=(S,S^\vee,R) \tag*{(93)}
\]
\[
\begin{aligned}
 &Z\subset A\quad\text{un \unsure{sous-gr. ou|sous-groupe} commutatif donné}\\
 &\text{d'où}\quad\deleted{Z}\ \Pi_Z\subset\mathrm{Bij}(S,S^\vee),\ \text{sur lequel}\\
 &Z\text{ opère librement (a priori }\Pi_Z\text{ peut être vide, mais :)}
\end{aligned}\tag*{(94)}
\]
\[
 \text{un sous-}Z\text{-torseur }\pi\text{ de }\Pi_A,\tag*{(95)}
\]
On pose
\[
 G=\mathrm{Centr}_A(Z). \tag*{(96)}
\]
\[
 \deleted{\unsure{G'|\mathcal G'}=\{g\in G\mid\uline{\text{centralisant }Z}.}
 \tag*{\deleted{\unsure{(97)|(98)}}}
\]

\uline{Lemme} Soit $f\in\deleted{\unsure{G}}G$, Conditions équivalentes

(i) $\exists\ \varphi_0\in\pi$ tel que $f(A_{\varphi_0})$ de \deleted{\unsure{A}} la forme $A_\varphi$
avec $\varphi\in\pi$

(i bis) Pour \uline{tout} $\varphi_0\in\pi$, on a $f(A_{\varphi_0})$ de la forme $A_\varphi$
pour $\varphi=f(\varphi_0)$ convenable -- i.e. $f(\pi)=\pi$
(transport de structure)

\deleted{(ii) $\exists f\in$}
\deleted{L'ens. de ces $f$ forment un sous}

Les $f$ satisfaisant ces conditions \deleted{sont} $\in G$,
et forment un sous-groupe \inserted{$\unsure{\mathbf G|G}$} de $G$, contenant
\deleted{les} groupes \deleted{\unsure{Z_2|\mathcal A}} $\mathbf G_0=\mathrm{Aut}(S,A_{\varphi\deleted{\unsure{0}}})$ \deleted{pour}
\deleted{\unsure{$\varphi\in$} toute $\varphi\in\pi$, qui \unsure{est aussi}}
\inserted{(indépendant du choix de $\varphi$) et $Z$.} C'est aussi le sous-groupe
du groupe des autom. de la structure de
multigraphe à homothéties $(S,S^\vee,R,Z,\pi)$
induisant l'identité sur $Z$. On a un
diagramme, où la ligne est exacte
\diagram{\begin{tikzcd}[ampersand replacement=\&,column sep=large]
{\text{(97)}} \& 1\arrow[r] \& \mathbf G_0\arrow[r] \& \mathbf G\arrow[r,"\delta"]\arrow[d] \& Z\\
\& \& \& Z\arrow[ur,swap,"{\lambda\mapsto\lambda^2}"] \&
\end{tikzcd}}
% archive sheet number: 220

\scanpage{222}
La ``condition d'isomorphie'' (ou première
condition d'homogénéité) dit que l'hom
$G\xrightarrow{\delta}Z$ est surjectif, i.e. que l'on a
une suite exacte courte
\diagram{\begin{tikzcd}[ampersand replacement=\&,column sep=large]
{\text{(98)}} \& 1\arrow[r,dash] \& G_0\arrow[r,dash] \& G\arrow[r,"\delta"] \& Z\arrow[r,dash] \& 1 \& {.}\\
\& \& \& Z\arrow[u,hook]\arrow[ur,swap,"{\lambda\mathrel{\unsure{\mapstochar\relbar|\mapsto}}\lambda^2}"] \& \& \&
\end{tikzcd}}

\deleted{J'aimerais maintenant, dans le contexte}

\deleted{Je suppose maintenant que la structure
de graphe \deleted{\unsure{envis.}} $(S,A_0)$, correspondant à un
$i_0\in\pi$, est triangulaire et même associée
à une structure de $2$-polyèdre \inserted{connexe}, et que
$G_0$ soit transitif sur l'ensemble des
repères. Je vais faire bon poids d'hypothèses,
quitte à raffiner par la suite, en supposant
le $2$-polyèdre $P_0$ orientable, soit $G_0^+$ le sous-groupe
d'indice $2$ de $G_0$ qui conserve l'orientation
-- on \unsure{soit|voit} que c'est le même pour $P_0$ et
pour les $P_i$ ($i\in\pi$). A vrai dire, ce qu'on
va utiliser, c'est la donnée d'un sous-groupe}
\[
 \deleted{G_0^+\subset G_0} \tag*{\deleted{(7)}}
\]
\deleted{qui ait les propriétés}
\[
 \left\{\begin{aligned}
 \deleted{\text{a) } G_0^+ \text{ d'indice }2\text{ dans }G_0,\text{ et invariant dans }G}\\
 \deleted{\text{b) } G_0^+\text{ est \uline{simplement} \uline{transitif} sur l'ens }\overline A_0}\\
 \deleted{\text{des arêtes orientées de }P_0}
 \end{aligned}\right. \tag*{\deleted{(8)}}
\]
% archive sheet number: 221

\scanpage{223}
\mspage{114}
$\overline{\underline{\unsure{\mathrm{X})|\mathrm{XI}|\mathrm{XI})}}}$ \uline{Récapitulation}

\deleted{\unsure{4. 2.|4.} Récapitulation}

\uline{1}. soit
% slanted
\margin{Graphes, prégraphes, jauges, multigraphes, multijauges\ldots}
\[
 \Gamma=(\overset{\text{sommets}}S,\overset{\text{cosommets}}{S^\vee},
 \overset{\text{incidence}}R\deleted{\unsure{\ni}}\subset S\times S^\vee) \tag*{(1)}
\]
un ``pré-graphe''. Deux applications
\[
 \varphi:S\longrightarrow S^\vee,\qquad \psi^{\deleted{\unsure{\vee}}}:S^{\deleted{\unsure{\vee}}}\longrightarrow S^\vee
 \tag*{(2)}
\]
sont dites ``associées'' (ou ``$R$-associées'') si on a
\[
 \forall s,t\in S\qquad (s,\psi(t))\in R \quad\Longleftrightarrow\quad
 \deleted{\varphi(s),\psi}(t,\varphi(s))\in R. \tag*{(3)}
\]
\deleted{\unsure{Si}} \deleted{\unsure{i.e.}} \deleted{$\unsure{t\varphi sRt}$}

C'est une relation symétrique en $\varphi,\psi$. Si $R$ est
séparée en $S^\vee$ [i.e.
$\begin{aligned} S^\vee&\longleftrightarrow\mathcal P(S)\\
s^\vee&\longmapsto{}^sR(s^\vee)\end{aligned}$ injective]
alors pour tout $\varphi:S\longrightarrow S^\vee$, il existe au plus une
$\psi:S\longrightarrow S^\vee$ associée -- \deleted{on dit alors que $\varphi$}
\inserted{en tous cas, on} appelle quasi-jauges de $\Gamma$ les $\varphi:S\longrightarrow S^\vee$ telles
qu'il existe une $\psi$ associée.
\deleted{D'où sur l'ens des}
\inserted{Sur l'ens des quasi-jauges, on a donc, dans le cas $\Gamma$ sép. en $S^\vee$, une involution
$(3')\ \varphi\longmapsto{}^s\varphi$.}
\deleted{quasi-jauges, il y a une structure de graphe,
\unsure{définie par} les ``arêtes'' sont les parties de la forme
$\{\varphi,\psi\}$, avec $\varphi$ et $\psi$ associées et sur l'ensemble de
\unsure{celles-ci}. Dans le cas $\Gamma$ séparé en $S^\vee$, une $\psi$ est une}

On appelle \uline{jauge} une $\varphi$ associée à elle $=$,
\unsure{i.e.|i.e. est} telle que
\[
 \forall(s,t)\in S\times S,\qquad (s,\varphi(t))\in R
 \ \Longrightarrow\ \deleted{\unsure{\varphi(s)}}(t,\varphi(s))\in R \tag*{(4)}
\]
ou encore, dans le cas où $\Gamma$ est séparé en $S^\vee$
\[
 \varphi={}^s\varphi .\tag*{(5)}
\]

Supposons que \deleted{l'on ait} $\varphi$ soit une bijection
\inserted{(on dit alors que $\varphi$ est une \uline{jauge stricte})}
\[
 \varphi:S\xrightarrow{\sim}S^\vee\qquad\text{d'où }\varphi'^{\deleted{\unsure{\vee}}}=\varphi^{-1}:S^\vee\xrightarrow{\sim}S
 \tag*{(6)}
\]
et considérons \deleted{l'application} \inserted{l'isom.} de $\underline{(\mathrm{Ens})\times(\mathrm{Ens})}$
\[
 (S,S^\vee)\xrightarrow[(\varphi,\varphi'\overset{\mathrm{déf}}=\deleted{\unsure{\vee}}\varphi^{-1})]{}
 (S'\overset{\mathrm{déf}}=S^\vee,S'{}^\vee\overset{\mathrm{déf}}=S)
 \tag*{(7)}
\]
en sorte que \deleted{\unsure{cette|ce} couple de bijections transforme}
% archive sheet number: 222

\scanpage{224}
% slanted
\margin{Cela marche dès que l'on a $\deleted{\varphi}: S\underset{\varphi'}{\overset{\varphi}{\rightleftarrows}}S^\vee$
avec $\varphi'\circ\varphi=\mathrm{id}_{\unsure{S^\vee|S}}$, alors
$(\varphi,\varphi')$ est \uline{compatible} avec $R,R'$ ssi $\varphi$ est une jauge}
\[
 \varphi\times\varphi'\text{ transforme }R\text{ en }R'\overset{\mathrm{déf}}={}^sR
 \qquad\text{ssi }\varphi\text{ est une jauge.}\tag*{(8)}
\]
Ce qui donne une correspondance biunivoque \deleted{\unsure{une}}
\[
 \underset{\text{strictes}}{\mathrm{Jauges}}(\Gamma)\simeq
 \begin{aligned}[t]&\text{Ens des isos }(\varphi,\varphi'):\Gamma\longrightarrow\Gamma'\\
 &\text{tels que }\varphi\deleted{\unsure{\varphi}},\varphi'\text{ soient inverses}\\
 &\text{l'un de l'autre, où }\Gamma'\text{ est le}\\
 &\text{prégraphe }(S^\vee,S,{}^sR)\text{ dual de }\Gamma.
 \end{aligned}\tag*{(9)}
\]

Si $\varphi$ est une jauge de $\Gamma$, alors
\[
 \Gamma_\varphi=(S,A_\varphi),\quad\text{où }\ 
 \underset{\substack{\Vert\\\{\{s,t\}\mid\{s,\varphi(t)\}\in R\}}}{A_\varphi}\subset\mathcal P(S)
 \ ,\quad\begin{aligned}[t]&\text{est un graphe}\\&\text{(avec des arêtes}\\
 &\text{qui peuvent être}\\&\text{ponctuelles)}
 \end{aligned}\tag*{(10)}
\]
la relation symétrique associée à cette str. de
graphe étant
\[
 R_\varphi=(\mathrm{id}_S\times\varphi)^{-1}(R)=(\varphi\times\mathrm{id}_S)^{-1}(R)
 =\{(s,t)\in S\times S\mid\{s,\varphi(t)\}\in R\}\tag*{(11)}
\]
-- c'est le \uline{graphe associé à} $\varphi$, \deleted{\unsure{noté} $\Gamma_\varphi$ Si $R$ est sép}
\inserted{On trouve ainsi} \deleted{sép} une application
\[
 \mathrm{Jauges}(\Gamma)\longrightarrow
 \begin{aligned}[t]&\text{ens des structures}\\&\text{de graphes }\gamma\text{ sur }S\\
 &\text{telles que l'ens des étoiles }\mathcal E_\gamma\\
 &\text{correspondants soit }\deleted{\text{égal à}}\inserted{\text{contenu dans}}\\
 &\mathcal E=\mathrm{Im}(S^\vee\longrightarrow\mathcal P(S))\subset\mathcal P(S)
 \end{aligned}\tag*{(12)}
\]
% slanted
\margin{$\mathcal E\inserted{\subset\mathcal P(S)}$ s'appelle l'ens des \uline{étoiles} du prégraphe, on définit de $=$
$\mathcal E^\vee\subset\mathcal P(S^\vee)$, l'ens des coétoiles (une étoile est un ens de sommets, une coétoile un ens. de cosommets\ldots)}
Cette application étant \uline{bijective} si $\Gamma$ est séparée en $S^\vee$,
i.e. si on peut identifier $S^\vee$ à \deleted{\unsure{l'}} un sous-ens de $\mathcal P(S)$
\[
 S^\vee\hookrightarrow\mathcal P(S).\tag*{(13)}
\]
\deleted{Dans ce cas, \unsure{on}} On trouve une \unsure{réflection|restriction|application} \unsure{en dessous|évidente|en deçà}
\[
 \mathrm{Jauges\ surj}(\Gamma)\longrightarrow
 \begin{aligned}[t]&\{\text{ens des str. de}\\&\text{graphes }\inserted{\gamma}\text{ sur }S\\
 &\text{telles que }\mathcal E_\gamma=\mathcal E\\
 &[\text{i.e. ayant un ens}\\&\text{d'étoiles donné}
 \end{aligned}\tag*{(14)}
\]
On trouve aussi
% archive sheet number: 223

\scanpage{225}
\mspage{115}
(15) \unsure{Équivalence|Équivalences} de groupoïdes
\diagram{\begin{tikzcd}[ampersand replacement=\&,column sep=small]
{\begin{gathered}\text{groupoïde des graphes}\\(S,A_0\subset\mathcal P_{\{1,2\}}\deleted{\unsure{6}}(S))\end{gathered}}
\arrow[r,"\approx"] \&
{\begin{gathered}\text{groupoïde des prégraphes}\\\Gamma=(S,S^\vee,R)\text{ séparés en }S^\vee\\
\text{munis d'une jauge }\inserted{\text{surjective}}\ \varphi\end{gathered}}
\end{tikzcd}}
induisant à son tour

(16) Equivalence de groupoïdes \deleted{\unsure{$\longrightarrow$}}
\diagram{\begin{tikzcd}[ampersand replacement=\&,column sep=small]
{\begin{gathered}\text{groupoïde des graphes}\\(S,A_0)\text{ séparés (tout sommet}\\
\text{déterminé par son étoile)}\end{gathered}}\arrow[r,"\approx"] \&
{\begin{gathered}\text{groupoïde des prégraphes}\\\text{biséparés}\inserted{\Gamma},\text{ munis d'}\\
\text{une jauge stricte}\\\text{(i.e. bijective) }\varphi,\\(\text{ou encore (9) d'un}\\
\text{isom }(\varphi,\varphi'):\Gamma\xrightarrow{\sim}\Gamma'\\
\text{tel que }\varphi,\varphi'\text{ inverses l'}\\\text{un de l'autre, où }\Gamma'\\
\text{prégraphe dual) .}\end{gathered}}
\end{tikzcd}}

Soit maintenant
\[
\deleted{\unsure{\lambda}}\deleted{\unsure{\in}} A=\mathrm{Aut}(\Gamma)=
\left\{\inserted{\lambda=}(\lambda_S,\lambda_{S^\vee})\in
\unsure{\mathfrak S|\mathfrak G}_S\times\mathfrak S_{S^\vee}\ \middle|\
\begin{aligned}&\forall(s,s^\vee)\in S\times S^\vee\\
&\deleted{\unsure{\lambda_Ss}}\\
&(s,s^\vee)\in R\Longleftrightarrow\\
&(\lambda_Ss,\lambda_{S^\vee}s^\vee)\in R
\end{aligned}\right\} \tag*{(17)}
\]
\deleted{et \unsure{et} soit $\varphi$ une jauge, d'où une structure de graphe
$\gamma_\varphi=(S,A_\varphi)$. Alors}
\[
\begin{aligned}
&\deleted{\text{Conditions équivalentes}}\\
&\deleted{\text{a) }\lambda_{S^\vee}\varphi=\varphi\lambda_S^{-1}}\\
&\deleted{\text{b) }\lambda_S\text{ est un}}
\end{aligned}\tag*{\deleted{(18)}}
\]
Notons que
\[
\begin{aligned}
&\text{Si }\Gamma\text{ est}\qquad \lambda\longmapsto\lambda_S
  \qquad\text{induit un }\deleted{\text{biject}}\inserted{\text{isom.}}\text{ du groupe}\\
&\text{séparé en }S^\vee\qquad A\overset{\unsure{\supset}}{\longrightarrow}\mathfrak S_S
  \qquad A\text{ avec le groupe des}\\
&\text{i.e. }S^\vee\hookrightarrow\mathcal P(S)
  \hspace{10em}\text{permutations }\deleted{\unsure{\lambda}}f\text{ de }S\\
&\text{alors}\hspace{14em}\text{telles que }\underbrace{f(\mathcal E)}_{\text{tr. de str.}}=\mathcal E\quad\text{i.e.}\\
&\text{où }\mathcal E=\mathrm{Im}(S^\vee\subset\mathcal P(S))\qquad
 \text{qui transforment étoile en}\\
&\hspace{15em}\text{étoile,}
\end{aligned}\tag*{(18)}
\]
\deleted{\unsure{Soit|Soient}} Soit maintenant $\varphi$ une jauge surjective de $\Gamma$ \inserted{séparé en $S^\vee$},
\deleted{définissant donc} \inserted{\unsure{fourn. donc|fournissant|correspondant} \unsure{:|à}} une structure de graphe
\[
 \gamma=\gamma_\varphi=(S,A_\varphi)\quad\text{avec }
 \underbrace{\mathcal E_\gamma}_{\substack{\text{ens des}\\\text{étoiles de }\gamma}}=\mathcal E
 \tag*{(19)}
\]
\deleted{On a alors une} et soit $f\in\mathfrak S_S$. On a alors
% archive sheet number: 224

\scanpage{226}
(20) On a \unsure{équivalence|équivalences}
\[
\begin{aligned}
&f\text{ est une \uline{homothétie} pour }\gamma_\varphi
 &&\Longleftrightarrow\ \exists\lambda\in A,\text{ avec }f=\lambda_S,\\
&\text{i.e. }\forall(s,t)\in S\times S,\quad\{s,t\}\in A_\varphi\Longleftrightarrow
 \{f^{-1}(s),f(t)\}\in A_\varphi &&\text{et}\\
&\text{i.e. }\forall t\in S,\quad f(\mathrm{Et}_{A_\varphi}(t))=\mathrm{Et}_{A_\varphi}(f^{-1}(s))
 &&\lambda_{S^\vee}\varphi=\varphi\lambda_S^{-1}
\end{aligned}
\]
On trouve ainsi
\[
 \begin{aligned}
 &\text{Correspondance} &\text{homothéties} &&\text{Automorphismes }\lambda\in A\\
 &\text{biunivoque} &\text{du graphe} &\simeq&\text{du prégraphe }\Gamma,\text{ tels que}\\
 &\text{entre} &\gamma_\varphi=(S,A_\varphi) &&\lambda_{S^\vee}\varphi=\varphi\lambda_S^{-1}.
 \end{aligned}\tag*{(21)}
\]

On dit qu'une jauge $\varphi$ \deleted{\unsure{est une}} (sans supposer \unsure{néc.t|ici} $\Gamma$
séparé en $S^\vee$ ni $\varphi$ surjective) est compatible avec
l'autom $\lambda$, ou est une $\lambda$-jauge, si elle satisfait
\[
 \lambda_{S^\vee}\varphi=\varphi\lambda_S^{-1}\qquad\text{i.e. }\lambda_{S^\vee}\varphi\lambda_S=\varphi
 \quad\deleted{\unsure{\varphi t}}\tag*{(22)}
\]
% slanted
\margin{\uline{NB} Si $\varphi$ bijectif, la condition \unsure{(23)|(22)} s'écrit aussi
$(22')\ \varphi(\lambda)=\lambda_{S^\vee}^{-1}$ ou
$(22'')\ \varphi^{-1}(\lambda_{S^\vee})=\lambda_S^{-1}$}
(de $=$, si $Z\subset A$, on définit les $Z$-jauges
comme les jauges satisfaisant (22) $\forall\lambda\in Z$).
Ainsi, Si $\varphi$ est une $\lambda$-\unsure{gauche|jauche|jauge}, alors \deleted{\unsure{\varphi}} $\lambda_S$ est
une homothétie pour $\gamma_\varphi$, et inversement dans
le cas où $\Gamma$ est séparé en $S^\vee$.

\deleted{Supposons} On a
\[
\begin{aligned}
&\text{Si }\varphi\text{ est une }\deleted{\text{jau}}\{\lambda,\mu\}\text{-jauge }\inserted{\text{bijective}},
 \text{ où }\lambda,\mu\in A,\\
&\text{alors }\varphi\text{ est une }\lambda\mu\text{-jauge sss }\lambda\mu=\mu\lambda
\end{aligned}\tag*{(23)}
\]
ce qui implique que \deleted{pour} \inserted{deux homothéties d'}un graphe séparé,
\deleted{\unsure{deux}} \inserted{leur composé} \inserted{est une} homothétie sss elles commutent.
Si $Z$ est un sss-groupe de $A$, tel qu'il existe une $Z$-jauge
bijective, alors $Z$ doit être commutatif.

Donnons nous maintenant \deleted{\unsure{un}}
\[
 Z\subset A=\mathrm{Aut}(\Gamma)\qquad\text{un sous-groupe \uline{commutatif}}
 \tag*{(24)}
\]
et considérons \deleted{une $Z$-jauge $\varphi_0$.}
% archive sheet number: 225

\scanpage{227}
\mspage{116}
\[
 \Pi_Z^{\wedge}=\text{ens des }Z\text{-jauges }\deleted{\text{bijectives}}\text{ de }\Gamma
 \tag*{(25)}
\]
alors $Z$ opère sur $\Pi_Z^\wedge$ par
\[
 T_\lambda\varphi\overset{\mathrm{déf}}=\lambda_{S^\vee}\circ\varphi=\varphi\circ\lambda_S^{-1}
 \tag*{(25)}
\]
-- l'opération correspondante sur l'ens des graphes
associés aux $\varphi$ étant donnée par
\[
 T_\lambda A_\varphi=\{\{s,\lambda_S(t)\}\in\mathcal P(S)\mid s,t\ \text{t.q. }\{s,t\}\in A_\varphi\}.
 \tag*{(27)}
\]
Si $\varphi$ est surjective, alors pour $\varphi$ fixé, la connaissance
de $T_\lambda\varphi=\lambda_{S^\vee}\circ\varphi$ implique celle de $\lambda_{S^\vee}$, si $\varphi$ est injective
elle implique celles de $\lambda_S$, donc si $\varphi$ est bijective,
la connaissance de $T_\lambda\varphi$ implique celle de $\lambda$, i.e.
\deleted{\unsure{l'opérat|l'opération}}

(27) sur l'ens \inserted{$\Pi_Z$} des $Z$-jauges \uline{bijectives}, $Z$ opère
\uline{librement} -- donc sur \deleted{un graphe sé} \inserted{l'ens des}
structures de graphes \uline{séparés} sur un ens $S$,
admettant un sous-groupe comm. $Z_S\subset\unsure{\mathfrak S|\mathfrak G}_S$
comme groupe d'homothétie, $Z$ opère
librement.

Notons
\[
 G=\mathrm{Cent}_A(Z)\deleted{\unsure{s}},\text{ d'où }\quad Z\subset G\subset A
 \tag*{(28)}
\]
On trouve
\[
\left\{\begin{aligned}
 &\text{Soit }\varphi_0\in\Pi_Z\text{ une }Z\text{-jauge stricte }\inserted{\text{(i.e. bij)}}\text{ de }\Gamma,
 \ \deleted{\text{alors}}\inserted{\text{soit}}\\
 &\deleted{\unsure{G_{\varphi_0}|\mathcal G_{\varphi_0}}}\quad G_{\varphi_0}=\{g\in G\mid g_{S^\vee}\varphi=\varphi_{\unsure{0|\vphantom{0}}}g_S^{-1}
 \quad\text{i.e. }\varphi_0\text{ est une }g\text{-jauge}\}\\
 &(\text{\uline{NB} }G_{\varphi_0}\text{ n'est pas néc. un sous-groupe de }G)\\
 &\text{alors}\qquad\deleted{\text{alors}}\quad\deleted{\unsure{(30)}}\\
 &g\longmapsto\deleted{\unsure{g}}T_g\varphi\overset{\mathrm{déf}}=g_{S^\vee}\varphi_0=\varphi_0g_S^{-1}\\
 &\text{est une bijection}\\
 &G_{\varphi_0}\xrightarrow{\sim}\Pi_Z
\end{aligned}\right.\tag*{(29)}
\]
Mais on fera attention que par (23), $G_{\varphi_0}$ est un ss-groupe
de $G$ sss il est formé d'éléments $2:2$ permutables --
auquel cas c'est un sous-groupe commutatif, \deleted{\unsure{tel que}} \inserted{Dans ce}
% archive sheet number: 226

\scanpage{228}
cas et dans ce cas seulement, on peut considérer
que $\widehat Z=G_{\varphi_{\unsure{0|\varphi}}}$ opère sur $\Pi_Z$ (et non seulement sur $\varphi_0$),
-- et alors $\Pi_Z=\Pi_{\widehat Z}$ est un torseur sous $\widehat Z$.

Nous nous intéressons, pour $(\Gamma,Z\subset A_\Gamma=\mathrm{Aut}(\Gamma))$
donné, avec $\Gamma=(S,S^\vee,R)$, $Z$ ss-groupe commutatif,
\deleted{aux} \inserted{à} la donnée d'un sous-torseur $\pi$ de $\Pi_Z$,
sous l'opération de $\varphi$. Dans le cas d'une
$\Gamma$ biséparée, \unsure{cel|cela} signifie qu'on s'intéresse aux
``multigraphes séparés à groupe d'homothéties $Z$.''
\deleted{\unsure{On}}
\[
 \underset{\substack{\text{sous}\\Z\text{-torseur}}}{\pi}\subset\Pi_Z=\text{ens des } Z\text{-jauges bijectives}.
 \tag*{(30)}
\]
On trouve alors des bijections canoniques
\[
 S^\vee\simeq\pi\wedge_Z(S!),\qquad S\simeq\pi^\vee\wedge_Z(S^\vee!)
 \tag*{(31)}
\]
dont la \unsure{deuxième|dernière} peut être considérée, soit comme
déduite de la première, soit comme étant
l'isomorphisme du $=$ type que la première,
relatif à la situation duale
\[
 \left(\underset{\substack{\Vert\\(S^\vee,S,{}^sR)}}{\Gamma'},
 \underset{\substack{\Vert\\ {}^sZ=\{\{\lambda_{S^\vee},\lambda_S\}\mid\lambda\in Z\}}}{Z'},
 \underset{=\{\varphi^{-1}\mid\varphi\in\pi\}}{\pi^{\deleted{\unsure{\vee}}'}\ \deleted{\unsure{Z'}}}\right)
 \tag*{(32)}
\]
en notant que
\[
 \varphi\longmapsto\varphi^{-1}\qquad \pi\longrightarrow\pi'
 \tag*{\deleted{\unsure{(33)}}}
\]
\deleted{\unsure{est un}} \inserted{satisfait à}
\[
 (T_\lambda\varphi)^{-1}=T'_{\lambda^{-1}}(\varphi^{-1})
\]
et définit donc un isom can. de $Z$-torseurs
\[
\begin{aligned}
 \pi'\simeq\pi^{[-1]}\overset{\mathrm{déf}}={}&\pi\text{ avec opn de }Z\text{ opposée de}\\
 &\text{l'opération donnée } T_\lambda,\text{ i.e.}\\
 &\lambda\text{ opérant par }T_{\lambda^{-1}}\ (\overset{\mathrm{déf}}=T'_\lambda).
\end{aligned}\tag*{(33)}
\]
% archive sheet number: 227

\scanpage{229}
\mspage{117}
\deleted{L'opération de $Z$}

Considérons
\[
 S^{\unsure{\leftrightharpoons|\natural|\leftrightarrows|\flat}}=S/Z_S,\qquad S^{\vee\leftrightharpoons}=S^\vee/Z_{S^\vee}
 \tag*{(34)}
\]
alors les $\varphi\in\pi$ définissent par passage aux quotients
une $=$ application bijective
\[
 \varphi^\leftrightharpoons:S^\leftrightharpoons\xrightarrow{\sim}S^{\vee\leftrightharpoons}
 \tag*{(35)}
\]
par quoi on identifie $S^\leftrightharpoons,S^{\vee\leftrightharpoons}$ à un $=$ ensemble $B$.
\deleted{On pose de plus} Les orbites \inserted{\unsure{$Z$, ou de|$Z$-orbites}} de $Z_S$ dans $S$ s'appellent \uline{rayons}, \deleted{de \unsure{$B$},}
celles dans $S^\vee$ s'appellent corayons -- donc la donnée
de $\pi$ nous définit une bijection entre \inserted{ens. des} rayons et
\inserted{ens. des} corayons. On posera
\[
 R^\leftrightharpoons=R/Z \tag*{(36)}
\]
et \deleted{l'application $\subset$} \inserted{l'inclusion} $R\subset S\times S^\vee$ définit par passage
aux quotients
\[
 R^\leftrightharpoons\longrightarrow S^\leftrightharpoons\times S^{\vee\leftrightharpoons}\tag*{(37)}
\]
qui n'est plus injective en général, \deleted{En fait,} \inserted{mais qui l'est souvent.}
\deleted{L'opération de $Z$ sur $R$ \unsure{provient|procède|provenant} d'une opération de
$Z\times Z$ sur $S\times S^\vee$}
\deleted{sur} Supposons pour simplifier que $Z$ opère librement
sur $S$ \deleted{et} $S^\vee$ \inserted{(sur $S^\vee$ et sur la \unsure{relation})}
(une hypoth. implique l'autre d'ailleurs
quand on en a une en plus un $\pi$\ldots).
% slanted
\margin{hypothèse vérifiée dans le cas biséparé ``homogène'' \deleted{\unsure{qu'on suppose} \illegible{2}}}
Alors les rayons et les corayons sont des torseurs sous $Z$.

Soient \deleted{$\unsure{D\in S}$,}
\[
 \underset{\substack{\rotatebox{90}{$\in$}\\S^\leftrightharpoons}}D\subset S,\quad
 \underset{\substack{\rotatebox{90}{$\in$}\\S^{\vee\leftrightharpoons}}}{D^\vee}\subset S^\vee
\]
\deleted{deux} rayon et corayon ``liés'' i.e. \deleted{formant} correspondant
à un $(s^\leftrightharpoons,s^{\vee\leftrightharpoons})\in S^\leftrightharpoons\times S^{\vee\leftrightharpoons}$ dans l'image de $R^\leftrightharpoons$,
i.e. tels que
\[
 \exists s_0\in D,\ s_0^\vee\in D^\vee\quad\text{tels que }(s_0,s_0^\vee)\in R
 \qquad(\text{NB }D=Zs_0,\ D^\vee=Zs_0^\vee) \tag*{(38)}
\]
Alors la \deleted{partie} \inserted{fibre} de $R^\leftrightharpoons$ \deleted{sur} \inserted{en} $(s^{\leftrightharpoons},s^{\vee\leftrightharpoons})$ est le quotient
% archive sheet number: 228

\scanpage{230}
par $Z$ de la partie de $R$, stable sous $Z$,
\deleted{\unsure{formée de|définie}} \inserted{définie} par
\diagram{\begin{tikzcd}[ampersand replacement=\&,row sep=huge]
{\text{(39)}} \& {\deleted{\unsure{Z\in|Z}}\ R(s^{\unsure{\leftrightharpoons|\natural|\leftrightarrows|\flat}},s^{\vee\leftrightharpoons})=
\left\{(s,s^\vee)\in R\ \middle|\ \begin{aligned}&s\text{ au dessus de }s^\leftrightharpoons\\
&s^\vee\text{ au dessus de }s^{\vee\leftrightharpoons}\end{aligned}\right\}}\\
\& { \unsure{E|\mathcal E}(s_0,s_0^\vee)\overset{\mathrm{déf}}=
\left\{(\lambda,\mu)\in Z\times Z\ \middle|\
\begin{aligned}&(\lambda s_0,\mu s_0^\vee)\in R\quad\text{i.e.}\\
&((\mu^{-1}\lambda)s_0,s_0^\vee)\in R\quad\text{i.e.}\\
&\mu^{-1}\lambda s_0\in\underbrace{E(s_0^\vee)}_{\substack{\text{étoile dans }S\\
\text{définie par }s_0^\vee}}\end{aligned}\right\}.}
\arrow[u,"\wr" description,"{(\lambda,\mu)\longmapsto(\lambda s_0,\mu s_0^\vee)}"']
\end{tikzcd}}
On trouve ainsi une bijection can.
\[
\begin{aligned}
 &\text{fibre de }R^\leftrightharpoons\text{ en }(s^\leftrightharpoons,s^{\vee\leftrightharpoons})\simeq
 \deleted{\text{sous-groupe de}}\deleted{\text{fini}}
 \inserted{Z(s_0,s_0^\vee)}\ \inserted{\text{ens. des }}z\in Z\\
 &\hspace{13em}\text{tels que }\ (zs_0,s_0^\vee)\in R\\
 &\hspace{13em}\text{i.e. }\ (s_0,z^{-1}s_0^\vee)\in R\\
 &\hspace{13em}\text{i.e. }\ zs_0\in E
\end{aligned}\tag*{(40)}
\]
Cette bijection en fait, \inserted{qui dépend \inserted{du choix} de $s_0,s_0^\vee$,} ne dépend pas de l'hypothèse
$(s_0\,s_0^\vee)\in R$ -- les deux mbs de (40) \unsure{pouvant|peuvent} être vides,
\deleted{ce qui} L'hypothèse $(s_0,s_0^\vee)\in R$ signifie simplement que
$Z_{s_0,s_0^\vee}\ni1$. Quand on change de $(s_0,s_0^\vee)$, le
sous-ens. en question est remplacé par un translaté.

\deleted{Dans le cas où $Z$-jauge stricte $\varphi_0$, on a l'équivalence}
\inserted{On a l'équivalence [si $Z$ opère librement sur $S,S^\vee$]}
\[
\begin{aligned}
&\deleted{\text{équivalence}}\quad R^\leftrightharpoons\longrightarrow Z^\leftrightharpoons\times Z'{}^\leftrightharpoons
 \text{ injective}\ \Longleftrightarrow\
 \deleted{\unsure{E}}\inserted{\forall} E\in\mathcal E,\ s_0\in E,\\
&\hspace{12em}\lambda\in Z,\ \inserted{\text{on a}}\ \lambdas_0\in E\Longrightarrow\\
&\hspace{12em}\deleted{\lambda s_0\notin E}\quad\lambda=1\\
&\underset{}{\overset{\substack{\text{en termes d'une}\\Z\text{-jauge \uline{stricte}}}}{\Longleftrightarrow}}
 \ \forall(s,t)\in S\times S,\ \inserted{\text{on a }}\{s,t\}\in A_0\
 \deleted{\Longrightarrow}\ ,\ \{s,\lambda t\}\in A_0\Rightarrow\lambda=1\\
&\Longleftrightarrow\ \forall\lambda\in Z,\ \lambda\ne1,\ \text{on a }A_0\cap A_\lambda=\{\unsure{\emptyset}\}\\
&\deleted{\Longleftrightarrow\ \forall\lambda\in Z,\ \unsure{\lambda\ne1},\ A_0\cap A_\lambda=\emptyset
 \quad\text{si on donne une }Z\text{-jauge }\varphi_0,\ \lambda\in Z,\ \lambda\ne1}\\
&\deleted{\Longleftrightarrow\ \{\forall\lambda,\lambda'\in Z,\ \lambda\ne\lambda'\quad
 \text{sont disjointes}\quad A_\lambda\cap A_{\lambda'}=\emptyset}
\end{aligned}\tag*{(41)}
\]
Cette condition signifie que
si $D$ et $D^\vee$ sont \deleted{deux} \inserted{un} rayon et corayon incidents,
alors pour tout $s\in D$, $\exists!s^\vee\in D^\vee$ tel que $s,s^\vee\in D^\vee$.
Si $s$ remplacé par $\lambda s$, $s^\vee$ remplacé par $\lambda s^\vee$. Donc
\deleted{Si (41) \unsure{est tel}} \inserted{(41) équivaut à}
\[
\begin{aligned}
&\text{Pour deux rayons }\deleted{\unsure{D,D^\vee}}\inserted{D\in S^\leftrightharpoons,D^\vee\in S^{\vee\leftrightharpoons}}\text{ incidents,}\\
&\text{la relation }\deleted{\text{entre}}\ D\text{ et }D^\vee\qquad (s,s^\vee)\in R\text{ est}\\
&\text{une relation bijective, établissant \unsure{une isom|un isom|un isom.} de}\\
&Z\text{-torseurs}\\
&\hspace{6em}D\simeq D^\vee
\end{aligned}\tag*{$(41')$}
\]
% archive sheet number: 229

\scanpage{231}
\mspage{118}
L'interprétation en termes d'une \deleted{\unsure{$\varphi_0$}} $Z$-jauge
stricte, i.e. d'une structure de graphe $(S,A_0)$ ayant
$Z$ comme groupe d'homothétie, est un peu
différente

$(41'')$ (41) signifie que pour deux rayons $D,D'\in S^{\unsure{\leftrightharpoons|\natural|\leftrightarrows|\flat}}$
incidents (i.e. tels qu'il existe $(s,s')\in D\times D'$
avec $(s,s')\in A_0$), \deleted{\unsure{$\exists$}} $\forall s\in D$, $\exists!s'\in D'$ adjacent
à $\unsure{s_2|s,|s}$ -- cette bijection $D\simeq D'$ est un \uline{anti}-isomorphisme pour les structures de torseur,
donc établit un iso.
\[
 D\wedge_ZD'\simeq 1_Z.
\]

Enfin, en termes de la donnée d'un $\pi$, sous $Z$-torseur de $\Pi_Z$, on a une formulation un peu
différente encore. Introduisons
\[
 A=A_{\deleted{\unsure{\varphi}}\pi}=\bigcup_{i\in\pi}A_i\quad\subset\quad\mathcal P_{\{1,2\}}(S)
 \tag*{(42)}
\]
qui définit une structure de graphe sur $S$, on
a alors

\deleted{(43) Pour deux rayons $D,D'$ de $S$, \inserted{corr. à $s^\leftrightharpoons\in S^\leftrightharpoons,s'{}^\leftrightharpoons\in S'{}^\leftrightharpoons$,}
sont ``$\pi$-incidents'' i.e. \deleted{\unsure{soit $s^\leftrightharpoons$ et $s'$ sont}} \inserted{$s'{}^\leftrightharpoons=\varphi^\leftrightharpoons(s^\leftrightharpoons)$}
(35), \deleted{de $s^\leftrightharpoons$ par} ou encore $s^\leftrightharpoons=\varphi^\leftrightharpoons(s'{}^\leftrightharpoons)$,
sss $\exists s\in D,s'\in D'$ tels que $\{s,s'\}\in A$, ou encore sss $\forall s\in D,s'\in D'$,
\unsure{on a} $\{s,s'$}

(43) Soient \deleted{\unsure{$B,B'$}} $D,D'\subset S$, \deleted{des} \inserted{deux} rayons, correspondant
à $s^\leftrightharpoons,s'{}^\leftrightharpoons\in S^\leftrightharpoons$, conditions équivalentes
\[
\left\{\begin{aligned}
&\text{a) }s'{}^\leftrightharpoons=\varphi^\leftrightharpoons(s^\leftrightharpoons)\\
&\text{b) }s^\leftrightharpoons=\varphi^\leftrightharpoons(s'{}^\leftrightharpoons)\\
&\text{c) }\forall(s,s')\in D\times D',\ \text{on a}\quad\{s,s'\}\in A\qquad(42)\\
&\text{d) }\exists(s_0,s'_0)\in D\times D'\quad\text{avec }\{s_0,s'_0\}\in A
\end{aligned}\right.
\]
On dit alors que $D,D'$ sont \inserted{$\pi$-}\uline{incidents}, \deleted{ce qui} \inserted{condition}
qui ne dépend d'ailleurs que de la donnée de
$Z_S\subset\unsure{\mathfrak S|\mathfrak G}_S$ (permettant de définir $S^\leftrightharpoons$ et la notion de rayon) ;
% archive sheet number: 230

\scanpage{232}
et de la structure de graphe $(S\deleted{\unsure{_0}},A)$, et non
de la famille $\pi$ de sous-ens $A_i\subset A$, $i\in\pi$.
Ceci posé, on a

\deleted{(44) Les cond. équivalentes (41) (41') (41'') signifient
aussi que pour deux \inserted{$Z$-}rayons incidents,
$D,D'\subset S$, et pour tt $(s,s')\in D\times D'$, $\exists$ existe
un \uline{unique} $\deleted{\unsure{\lambda}}i=i_{s,s'}\in\pi$ tel que $\{s,s'\}\in A_i$, \deleted{On} \inserted{soit}
\deleted{de \unsure{int},} \deleted{a alors} On a $i_{\lambda s,s'}=i_{s,\lambda s'}$, et on trouve
ainsi un iso. canonique}
\[
 \deleted{D\wedge_ZD'\overset{\tau}{\simeq}\pi}\tag*{\deleted{(45)}}
\]
\deleted{\unsure{caractérisé|caractérisée} par $i_{s,s'}=\tau(s\wedge_Zs')$, i.e. par la \unsure{condition|égalité|relation}}
\[
 \deleted{\unsure{i_{s,s'}}}\tag*{\deleted{(46)}}
\]
(44) Les conditions \deleted{(41)} équivalentes (41), (41') (41'')
signifient aussi que les $A_i$ sont mutuellement
disjoints, i.e. \deleted{\unsure{on}} \inserted{on} a une application
canonique ``type''
% slanted
\margin{donc la structure de multigraphes $(A_i)_{i\in\pi}$ sur $S$ (avec un groupe d'homothéties $Z$)
est connue quand on connaît $S,A\deleted{\unsure{\bigcup_{i\in\pi}}}\subset\mathcal P_{\{1,2\}}(S)$,
sa partition en $A=\coprod_{i\in\pi}A_i$, et enfin le sous-groupe $Z\subset\unsure{\mathfrak S|\mathfrak G}_S$,
(avec la condition que $Z$ est formé d'homothéties pour les structures de graphe $(S,A_i)$, et que
\unsure{l'application} $\unsure{Z\ldots}$ transforme \unsure{canoniq.} les unes \unsure{dans les autres})}
\deleted{On trouve alors}
\[
 A\xrightarrow{\tau_A}\pi
 \tag*{\deleted{(45)}}
\]
\deleted{Si cette la} Si \inserted{$Z$} $D,D'$ sont \deleted{deux} rayons incidents,
\[
 (s,s')\longmapsto\tau_A(\{s,s'\})\qquad D\times D'\xrightarrow{\tau_{D,D'}}\pi
\]
satisfait
\[
 \deleted{\tau_A(\{\lambda s,s'\})}\quad
 \tau_{D,D'}(\lambda s,\mu s')=\deleted{\tau_{D,D'}(s,\lambda s')}
 \inserted{\lambda\mu\,\tau_{D,D'}(s,s')}, \tag*{(46)}
\]
et induit \deleted{donc} un isom. de $Z$-torseurs
\[
 D\wedge D'\xrightarrow{\sim}\pi .\tag*{(47)}
\]
% margin note written sideways
\margin{\deleted{Dans le cas}}
\uline{NB} La donnée d'un prégraphe $\Gamma$, avec groupe
$Z$ commutatif d'automorphismes -- jouant le
rôle d'homothéties virtuelles -- permet de
définir dans $S$ (disons, \unsure{et il suffit|et il faut|bien sûr} \deleted{\unsure{pour|plus}} $S^\vee$)
% archive sheet number: 231

\scanpage{233}
\mspage{119\deleted{\unsure{1|8}}}
deux types de sous-ensembles remarquables,
d'une part les \uline{étoiles} (indépendantes de la
donnée de $Z$). d'autre part les \uline{rayons} (ne
dépendant par contre que de la donnée de
$Z_{S}\subset\unsure{\mathfrak S|\mathfrak G}_S$). La relation \inserted{d'intersection} des rayons entre
eux est évidente (deux rayons distinct ont une
intersection vide), alors que celle de deux
étoiles peut-être plus complexe (il faudrait
l'étudier systématiquement \deleted{\unsure{?|$g$}}, \unsure{dans|donc} des cas où les
données structurales sont assez riches peut-être pour
faire une étude \deleted{\unsure{plus}} au moins suivie, p.ex.
le cas associé à un \deleted{\unsure{$E,S$}} $\Lambda$-module libre de
rang 2\ldots). La relation entre étoiles et rayons
est simple, dans le cas justement de l'hyp.
simplificatrice (41) et reformulations (41') (41''), (44).
\deleted{La considérant par} \deleted{Elle signifie en}
Notons \deleted{que ces conditions} \inserted{en effet} \deleted{\unsure{que, $s\in$}}
que si $s\in S.$ $s^\vee\in S^\vee$, et si $E$ \inserted{$\subset S$} est l'étoile définie par $s^\vee$,
alors
\[
 s\in E\quad\Longleftrightarrow\quad(s,s^\vee)\in R
\]
donc \deleted{l'hyp.\unsure{''}}

(48) \deleted{L'hyp.} \inserted{La condition} (41) signifie aussi que l'intersection dans $S$
d'une étoile \inserted{$E$} et d'un rayon \inserted{$D$} est, soit vide,
soit réduite à un point.

Dans le cas où on part d'une structure de graphe,
et où $E$ est l'étoile de $s\in E$, alors $E\cap D\ne\emptyset$
sss les rayons $D$ et \deleted{$\unsure{Z.s}$} $D'=Z.s$
sont incidents.
% archive sheet number: 232

\mspage{119 bis}
\uline{NB} La notion de prégraphe $\Gamma=(S,S^\vee,R)$ avec groupe $Z$ d'homothéties virtuelles
\[
 Z\subset\mathrm{Aut}(\Gamma),\qquad Z\text{ ss-gr. comm.}
\]

\scanpage{234}
sss les rayons
sont incidents.

\mspage{119 bis}
\uline{NB} La notion de prégraphe $\Gamma=(S,S^\vee,R)$ avec groupe $Z$ d'homothéties virtuelles
\[
 Z\subset\mathrm{Aut}(\Gamma),\qquad Z\text{ ss-gr. comm.} \tag*{(48.1)}
\]
satisfaisant les conditions
\[
 \left\{\begin{aligned}
 &\text{a) }Z\text{ opère librement sur }S\text{ et sur }S^\vee\\
 &\text{b) }\deleted{\unsure{Z}}\ \text{On a la condition (41) (ou (41')).}
 \end{aligned}\right.\tag*{(48.2)}
\]
peut s'expliciter \unsure{aussi|ainsi} \deleted{\unsure{aisém.}}. On peut identifier, à cause
de a), \deleted{\unsure{S}} $S$ et $S^\vee$ comme des $Z$-torseurs relatifs
sur deux ens $S^{\unsure{\leftrightharpoons|\natural|\leftrightarrows|\flat}}$ et $S^{\vee\leftrightharpoons}$, liés par une relation
\diagram{\begin{tikzcd}[ampersand replacement=\&,column sep=small]
{\text{(48.3)}} \& \& R^\leftrightharpoons\arrow[dl,swap,"{\unsure{p_1|p}}"]\arrow[dr,"p^\vee"] \&
 {\subset\underset{\text{ens des rayons}}{S^\leftrightharpoons}\times\underset{\text{ens des corayons}}{S^{\vee\leftrightharpoons}}} \&
 {\begin{gathered} (R^\leftrightharpoons\text{ la ``relation d'in-}\\\deleted{\text{relation d'incidence}}\\
 \text{cidence'' entre}\\\text{rayons et corayons)}\end{gathered}}\\
\& S^\leftrightharpoons \& \& S^{\vee\leftrightharpoons} \&
\end{tikzcd}}
et on a un isom de $Z$-torseurs
\[
 p^*(S/S^\leftrightharpoons)\simeq p^{\vee *}(S^\vee/S^{\vee\leftrightharpoons})
 \tag*{(48.4)}
\]
-- la relation $R$ entre $S$ et $S^\leftrightharpoons$ s'identifie à
$p^*(S/S^\leftrightharpoons)$ et $p^{\vee *}(S^\vee/S^{\vee\leftrightharpoons})$ identifiés par (44.4), s'envoyant
dans $S$ et dans $S^{\deleted{\unsure{*}}\vee}$ via les diagrammes commutatifs
\diagram{\begin{tikzcd}[ampersand replacement=\&,column sep=small]
\& S\arrow[d] \& {\deleted{\unsure{S^\leftrightharpoons}}}\arrow[d] \&
{\overbrace{p^*(S/S^\leftrightharpoons)\simeq p^{\vee *}(S^\vee/S)}^{R}}
\arrow[ll]\arrow[d,start anchor=south west,end anchor=north west]\arrow[d,start anchor=south east,end anchor=north east]\arrow[rr,<->]
\& {\deleted{\unsure{S}}} \& S^\vee\arrow[d]\\
{\text{(48.5)}} \& S^\leftrightharpoons \& {\deleted{\unsure{R^\leftrightharpoons|R'^{\leftrightharpoons}}}} \&
R^\leftrightharpoons\arrow[ll,"{\unsure{p_1|p_4|p_{\flat}}}"']\arrow[rr,swap,"p^\vee"] \& \& S^{\vee\leftrightharpoons}
\end{tikzcd}}

Ceci posé, \deleted{\unsure{on a} $\unsure{\longrightarrow}$} la donnée d'une \deleted{sous $Z$-torseur $\pi$ de $\Pi_Z$}
\inserted{$Z$-jauge stricte} peut \unsure{s'interpréter|s'interprêter} en termes
de
\[
\left\{\begin{aligned}
 &\text{a) une }\deleted{\unsure{\text{iso}}}\text{ bijection }S^\leftrightharpoons\xrightarrow{\varphi^\leftrightharpoons}S^{\vee\leftrightharpoons}
 \quad\text{qui soit}\\
 &\hspace{13em}\text{une }R^\leftrightharpoons\text{-jauge}\\
 &\deleted{\text{b) un }Z\text{-torseur }\pi}\\
 &\text{b) un }\deleted{\text{torseur}}\inserted{\text{isomorphisme}}\\
 &\deleted{\unsure{g}}(S/S^\leftrightharpoons)\wedge_{Z_{S^\leftrightharpoons}}(\varphi^{\leftrightharpoons *}(S^\vee/S^{\vee\leftrightharpoons}))
 \simeq 1_{\deleted{\unsure{Z}}S^\leftrightharpoons}
\end{aligned}\right.\tag*{(48.6)}
\]
\margin{(i.e. donnant lieu à une relation \uline{symétrique} sur
$S^\leftrightharpoons\times S^\leftrightharpoons$, \deleted{\unsure{soit} $Z\unsure{R}$} soit
$\unsure{R^\leftrightharpoons_{S^\leftrightharpoons}}$}
% slanted
\margin{\deleted{\unsure{la précisant} : la donnée d'une \unsure{bijection} qui soit \unsure{anticommutant} à l'opération de $Z$, i.e satisfasse $\varphi\circ\lambda=\lambda^{-1}\circ\varphi$}}
\deleted{ou autrement dit la donnée d'une \unsure{de}}
% archive sheet number: 233
